\chapter{Zadatak 1 --- od riječi do matrica}
\label{ch:z1-formulacija}

\section{Što se ovdje radi}

Zadatak je zapisan u ,,školskom`` obliku, s pojedinačnim koeficijentima. Sve
metode koje slijede --- KKT uvjeti, dualnost, \texttt{quadprog} --- traže matrični
zapis. Ovo poglavlje taj prijevod izvodi do kraja i utvrđuje notaciju koja
vrijedi do kraja zadatka.

\begin{predznanje}
\begin{itemize}[nosep,leftmargin=*]
  \item Kvadratna forma i čitanje $Q$ i $c$ iz funkcije cilja --- \S\ref{sec:kvadratna-forma}.
  \item Množenje matrice i vektora --- \S\ref{sec:mnozenje}.
  \item Zapis poluprostora $a\trans x \leq b$ --- \S\ref{sec:poluprostor}.
  \item Značenje simbola $\coloneqq$, $\leq$ po komponentama, $x\trans$ --- Poglavlje~\ref{ch:rjecnik}.
\end{itemize}
\end{predznanje}

\section{Standardni oblik koji koristi profesor}
\label{sec:z1-standardni-oblik}

\begin{izvor}{Predavanje 5 --- LP i QP, slajd 29/34 (PDF str.~33)}
\slajd[0.86]{p05-qp}

\textbf{Transkripcija.} \emph{Kvadratno programiranje (QP).} Zadano: matrice
$Q$, $A$, $G$ i vektori $c$, $b$, $h$.
\[
\begin{aligned}
\text{minimiziraj} \quad & \tfrac{1}{2}x\trans Q x + c\trans x \\
\text{uz ograničenja} \quad & Ax = b, \\
& Gx \leq h,
\end{aligned}
\]
gdje je $x \in \R^n$ vektor projektnih varijabli.
\begin{itemize}[nosep,leftmargin=*]
  \item Funkcija cilja je kvadratna dok su ograničenja linearne funkcije u
        projektnim varijablama.
  \item Matrica $Q$ mora biti pozitivno definitna (tada imamo (striktnu)
        konveksnost i nužno konačan minimum).
\end{itemize}

\textbf{Što ovo znači za naš zadatak:} ovo je definicija klase kojoj Zadatak 1
pripada. Ograničenja jednakosti nemamo, pa se dio $Ax = b$ ispušta. Ostaje
$\min \tfrac12 x\trans Qx + c\trans x$ uz $Gx \leq h$, a naš je posao samo
popuniti $Q$, $c$, $G$ i $h$ brojevima.
\end{izvor}

\section{Funkcija cilja u matričnom obliku}
\label{sec:z1-cilj-matricno}

Zadana funkcija cilja je
\begin{equation}
f(x) = x_1^2 + 2x_2^2 - 0{,}3\,x_1x_2 + 2x_1 - 3x_2 .
\label{eq:z1-f-skalarno}
\end{equation}

\begin{korak}[Očitavanje $Q$ i $c$]
Po pravilima iz \S\ref{sec:kvadratna-forma}:
\begin{enumerate}[leftmargin=*]
  \item Uz $x_1^2$ stoji koeficijent $1$, a u standardnom obliku je taj koeficijent
        $\tfrac{1}{2}Q_{11}$, pa je $Q_{11} = 2$.
  \item Uz $x_2^2$ stoji $2$, dakle $Q_{22} = 4$.
  \item Uz $x_1x_2$ stoji $-0{,}3$; u standardnom obliku mješoviti član iznosi
        $\tfrac{1}{2}\cdot 2Q_{12} = Q_{12}$, pa je $Q_{12} = Q_{21} = -0{,}3$.
  \item Linearni članovi daju $c_1 = 2$ i $c_2 = -3$.
\end{enumerate}
\end{korak}

\begin{formula}{Matrice funkcije cilja Zadatka 1}
\begin{equation}
Q = \begin{bmatrix} 2 & -0{,}3 \\ -0{,}3 & 4 \end{bmatrix}, \qquad
c = \begin{bmatrix} 2 \\ -3 \end{bmatrix},
\qquad
f(x) = \tfrac{1}{2}x\trans Q x + c\trans x .
\label{eq:z1-Qc}
\end{equation}
\end{formula}

Provjera unatrag:
$\tfrac12 x\trans Q x = \tfrac12(2x_1^2 - 0{,}6x_1x_2 + 4x_2^2) = x_1^2 - 0{,}3x_1x_2 + 2x_2^2$,
i $c\trans x = 2x_1 - 3x_2$. Zbroj je točno \eqref{eq:z1-f-skalarno}.

\section{Ograničenja u obliku $Gx - h \leq 0$}
\label{sec:z1-ogranicenja}

Uvrštavanjem zadanih koeficijenata $a_{11} = -1$, $a_{12} = -1$, $a_{21} = 1$,
$a_{22} = -1$ u (1b) i (1c) dobiva se
\begin{align}
-x_1 - x_2 &\leq -10, \label{eq:z1-ogr1-sirovo} \\
 x_1 - x_2 &\leq 3 . \label{eq:z1-ogr2-sirovo}
\end{align}

Prebacivanjem desne strane ulijevo, kako bi ograničenja imala oblik ,,funkcija
$\leq 0$``:
\begin{formula}{Ograničenja Zadatka 1}
\begin{align}
g_1(x) &= -x_1 - x_2 + 10 \leq 0, \label{eq:z1-g1} \\
g_2(x) &= x_1 - x_2 - 3 \leq 0, \label{eq:z1-g2}
\end{align}
odnosno matrično $Gx - h \leq 0$ uz
\begin{equation}
G = \begin{bmatrix} -1 & -1 \\ 1 & -1 \end{bmatrix}, \qquad
h = \begin{bmatrix} -10 \\ 3 \end{bmatrix}.
\label{eq:z1-Gh}
\end{equation}
\end{formula}

Provjera: $Gx - h = \begin{bsmallmatrix} -x_1 - x_2 \\ x_1 - x_2 \end{bsmallmatrix}
- \begin{bsmallmatrix} -10 \\ 3\end{bsmallmatrix}
= \begin{bsmallmatrix} -x_1 - x_2 + 10 \\ x_1 - x_2 - 3 \end{bsmallmatrix}
= \begin{bsmallmatrix} g_1(x) \\ g_2(x) \end{bsmallmatrix}$. Slaže se s
\eqref{eq:z1-g1}--\eqref{eq:z1-g2}.

\begin{zamka}[Zamka: predznak pri prebacivanju]
Ograničenje \eqref{eq:z1-ogr1-sirovo} ima negativnu desnu stranu, $-10$. Pri
prebacivanju se dobiva $+10$, ne $-10$:
\[
-x_1 - x_2 \leq -10
\;\Longleftrightarrow\;
-x_1 - x_2 + 10 \leq 0 .
\]
Ako se ovdje promaši predznak, cijeli dopustivi skup ispadne na krivoj strani
pravca, a KKT sustav onda ima rješenje s $\mu_1 < 0$, pa se odbaci točan slučaj.
Kontrola: uvrsti neku očito dopustivu točku, npr.\ $x = (7,\,6)$, i provjeri da
je $g_1 < 0$. Ovdje: $g_1(7,6) = -7 - 6 + 10 = -3 < 0$. \checkmark
\end{zamka}

\begin{zamka}[Zamka: $h$ nije ,,desna strana`` u uobičajenom smislu]
U zapisu $Gx - h \leq 0$ vektor $h$ sadrži \emph{izvorne} desne strane
$(-10,\; 3)$, a ne one prebačene. Predznak se pojavljuje tek u $Gx - h$. Ako
umjesto toga napišeš $h = (10,\,-3)$, matrični i skalarni zapis prestaju biti
isti problem. \texttt{quadprog} očekuje baš $Gx \leq h$, dakle $h = (-10,\;3)$.
\end{zamka}

\section{Dopustivi skup}
\label{sec:z1-dopustivi-skup}

\begin{pojam}{Dopustivi skup Zadatka 1}
\begin{equation}
\feas = \{\, x \in \R^2 \mid Gx - h \leq 0 \,\}
      = \{\, x \in \R^2 \mid x_1 + x_2 \geq 10 \ \text{ i } \ x_1 - x_2 \leq 3 \,\}.
\label{eq:z1-feas}
\end{equation}
\end{pojam}

Riječima: presjek dvaju zatvorenih poluprostora. Prvi traži da zbroj varijabli
bude barem $10$; drugi da $x_1$ ne prestigne $x_2$ za više od $3$. Presjek je
neograničeno klinasto područje čiji je vrh sjecište pravaca,
\begin{equation}
x_1 + x_2 = 10 \ \wedge \ x_1 - x_2 = 3
\;\Longrightarrow\;
x = \begin{bmatrix} 6{,}5 \\ 3{,}5 \end{bmatrix}.
\label{eq:z1-vrh}
\end{equation}
Ta točka će se pojaviti kao Slučaj 4 pri rješavanju KKT uvjeta.

\begin{predznanje}
Notacija koja od ovog trenutka vrijedi do kraja Zadatka 1:
\begin{itemize}[nosep,leftmargin=*]
  \item $x = [x_1,\, x_2]\trans \in \R^2$ --- projektne varijable;
  \item $Q$, $c$ iz \eqref{eq:z1-Qc}; $G$, $h$ iz \eqref{eq:z1-Gh};
  \item $g_1$, $g_2$ iz \eqref{eq:z1-g1}--\eqref{eq:z1-g2};
  \item $\mu = [\mu_1,\, \mu_2]\trans \geq 0$ --- Lagrangeovi množitelji;
  \item $\feas$ --- dopustivi skup \eqref{eq:z1-feas}.
\end{itemize}
\end{predznanje}

\begin{izgovor}
,,Problem sam najprije sveo na standardni oblik kvadratnog programa s
predavanja: minimiziraj jedna polovina iks transponirano ku iks plus ce
transponirano iks, uz ge iks minus ha manje ili jednako nuli. Matrica ku je dva,
minus nula tri, minus nula tri, četiri; vektor ce je dva, minus tri. Ograničenja
sam prebacio tako da desno stoji nula, pa je ge matrica minus jedan minus jedan,
jedan minus jedan, a ha vektor minus deset, tri.``
\end{izgovor}
